% Pursuing Stacks (cotes 134-2 to 134-8): archive cote 134-8, site page 33 (PDF page 34), typescript page 587.
% typescript. Draft compared: ps7.tex:1543-1608.
% AI-transcribed by Claude Opus 5.5 (claude-opus-5-5, adaptive thinking, effort high), Message Batches API, batch ps4 (msgbatch_012Aoc1NuskATPsyBeKSYK7J), 2026-11-12; not yet checked by eye.
\scanpage{134-8/33}
\mspage{587}
without necessarily grinding through a complete formulaire for the
five\inserted[typed]{main} operations we got so far (namely $Lf_{!}$,$f^{*}$,$Rf_{*}$, $\overset{L}{\otimes}$ , $R\underline{Hom}$ ). Not
later than two pages ago or so, we were faced again with two visibly
dual statements, one about $Lf_{!}$ , the other about $Rf_{*}$ - and feeling
silly not to be able to merely deduce one from the other !

29.10.

139. Last night I still pondered a little about the $Lf_{!}$ operation,
and did some more reading in the notes of chapter V of about three
months ago, which had been getting a little distant in my mind. In
those notes a great deal of emphasis goes with the notions of integra-%
tors and cointegrators - as a matter of fact, \xout{it was} that whole
chapter sprung from an attempt to understand the meaning \xout{and role} of
certain ``standard complexes'' associated with standard test categories
such as \hand{$\Delta$}, which then led us to the notion\xout{s} of an integrator (via the
intermediate one of an ``abelianizator''). This in turn brought us to
\margin{\textcircled{$*_k$}}the $*_{k}$ - formalism, expressing most conveniently the relationship
between categories such as $A^{\wedge}_{k}$ and $(A^{o})^{\wedge}_{k}$ (k any commutative ring),
and it's various avatars. I was so pleased with this formalism and it'
s ``computational'' flavor, that in my enthousiasm I subsumed under it
the dual treatment of homology and cohomology for an\xout{category (with}
object A of (Cat) (with coefficients in a complex of \hand{$\mathcal{M}$}-valued
presheaves,\hand{$\mathcal{M}$} being an abelian category satisfying some mild conditi-%
ons), in section 108, as a particular case of the total derived
functors $F_{.}\overset{L}{*}_{k}L'_{.}$ and $R\mathrm{Hom}_{k}(L_{.},K^{\cdot})$ ( with values in the derived
category of \hand{$\mathcal{M}$} ). The main point here was using \uline{projective resolutions}
of the argument \hand{$L_{.},L'_{.}$} in $B^{\wedge}_{k}$ or $A^{\wedge}_{k}$ (where $B=A^{o}$), which in \xout{most cases}
the most important case was just the constant sheaf of rings $k_{B}$ or $k_{A}$
- rather than resolve the argument $F_{.}$ (projectively) or $K^{\cdot}$ (injective-%
ly), namely the coefficients for homology or cohomology\hand{.}\deleted{respectively.}
It was the enthousiasm of the adept of a game he just discovered -
I was going to try it out for the next step, namely relative homology
and cohomology $Lf_{!}$ and $Rf_{*}$ , with f any map in (Cat)\deleted{,} - but then I
got caught by the more fascinating schematization game. Coming back
now upon the rather routine matter of looking up a comprehensive
mutually dual treatment of $Lf_{!}$ and $Rf_{*}$ , it does'nt seem that the
formalism of integrators and cointegrators is going to be of much help.
To be more specific, in order to compute (or simply define) $Lf_{!}(F_{.})$
\xout{or $Rf_{*}$}

\margin{33}
